% Part: counterfactuals
% Chapter: minimal-change-semantics
% Section: antecedent-strengthening

\documentclass[../../../include/open-logic-section]{subfiles}

\begin{document}

\olfileid[hi]{cnt}{min}{agg}

\olsection{पूर्ववर्ती का प्रबलीकरण}

``पूर्ववर्ती का प्रबलीकरण'' अनुमान
$!A \lif !C \Entails (!A \land !B) \lif !C$ को कहते हैं। यह भौतिक सशर्त
के लिए मान्य, किंतु प्रतितथ्यात्मक के लिए अमान्य है। मान लें यह सत्य है कि
यदि मैं इस माचिस की तीली को रगड़ता, तो वह जलती। (इसका अर्थ है कि तीली या
डिबिया की रगड़ने वाली सतह में कोई खराबी नहीं, मैं तीली तोड़ूँगा नहीं,
इत्यादि।) किंतु यह सत्य नहीं है कि यदि मैं बाह्य अंतरिक्ष में इस तीली को
जलाता, तो वह जलती। अतः निम्न अनुमान अमान्य है:
\begin{quote}
  यदि माचिस की तीली रगड़ी जाती, तो वह जलती।

  अतः यदि माचिस की तीली बाह्य अंतरिक्ष में रगड़ी जाती, तो वह जलती।
\end{quote}

सशर्तों का लुईस--स्टालनेकर विवरण इसे समझाता है: वह निकटतम जगत जहाँ मैं तीली
जलाता हूँ और ऐसा बाह्य अंतरिक्ष में करता हूँ, उस निकटतम जगत की अपेक्षा
वास्तविक जगत से बहुत अधिक दूर है जहाँ मैं तीली जलाता हूँ। अतः यद्यपि बाद
वाले में तीली जलती है, पहले वाले में नहीं। और ऐसा ही होना चाहिए।

\begin{ex}
  गोलक अर्थविज्ञान इस अनुमान को अमान्य करता है; अर्थात हमारे पास
  $p \cif r \Entails/ (p \land q) \cif r$ है। प्रतिरूप
  $\mModel{M} = \tuple{W, O, V}$ पर विचार करें, जहाँ
  $W = \{w, w_1, w_2\}$,
  $O_w = \{\{w\}, \{w, w_1\}, \{w, w_1, w_2\}\}$,
  $V(p) = \{w_1, w_2\}$, $V(q) = \{w_2\}$ और $V(r) = \{w_1\}$। एक
  $p$-स्वीकारी गोलक $S = \{w, w_1\}$ है और उसके सभी जगतों पर
  $p \lif r$ सत्य है, अतः $\mSat{M}{p \cif r}[w]$। एक
  $(p \land q)$-स्वीकारी गोलक $S' = \{w, w_1, w_2\}$ भी है, किंतु
  $\mSat/{M}{(p \land q) \lif r}[w_2]$; अतः
  $\mSat/{M}{(p \land q) \cif r}[w]$
  (\olref{fig:antecedent-strengthening} देखें)।

\begin{figure}
\begin{center}
\begin{tikzpicture}[layerwidth=1.5,scale=.6]\tiny
  \spheresystem{3}
  \propositionintersect{-10}{3}{30}{5}
  \begin{scope}
    \clip plot[smooth,tension=1] coordinates
          {(0,4.5) (30:.95) (4.5,-3.5)} -- (4.5,4.5) ;
    \spherefill{2}
  \end{scope}
  \draw[proposition] plot[dashed,smooth,tension=1] coordinates
          {(0,4.5) (30:.95) (4.5,-3.5)} ;
  \proposition{30}{2}{30}{5}
  \path (0,0) node[world] {$w$};
  \spherepos{30}{2}{node[world] {$w_1$}}
  \spherepos{-10}{3}{node[world] {$w_2$}}
  \spherepos{-10}{4}{node {$q$}}
  \spherepos{30}{4}{node {$r$}}
  \draw (1,4.5) node[below] {$p$};
\end{tikzpicture}
\caption{पूर्ववर्ती के प्रबलीकरण का प्रतिदृष्टांत}
\ollabel{fig:antecedent-strengthening}
\end{center}
\end{figure}
\end{ex}

\end{document}
